\section{Related Work}
\label{sec:related_work}

The evolution of automated phishing detection spans four major paradigms: heuristic/blacklist methods, content-based deep learning, URL/network-based machine learning, and Explainable AI (XAI) frameworks.

\subsection{Heuristic and Blacklist-Based Detection}
Early defenses relied on signature repositories and community blacklists \cite{sahoo2017malicious}. While services like PhishTank and Google Safe Browsing provide authoritative verdicts for cataloged threats, studies by Hannousse and Yahiouche \cite{hannousse2021towards} and Marchal et al. \cite{marchal2017off} show that blacklists suffer from significant discovery latency. Attackers evade blacklists via domain generation algorithms (DGAs), fast-flux DNS, and disposable top-level domains (TLDs), enabling zero-day attacks to execute before signatures propagate.

\subsection{Content-Based and Deep Learning Approaches}
To overcome blacklist delays, several studies proposed inspecting web page contents, DOM trees, and scripts. Opara et al. \cite{opara2020htmlphish} developed \textit{HTMLPhish}, employing CNN-LSTM networks to analyze raw HTML document flows. Mohammad et al. \cite{mohammad2014predicting} investigated self-structuring neural networks on page-level features. Adebowale et al. \cite{adebowale2023intelligent} integrated convolutional and recurrent architectures on website features, while Al-Alyan and Al-Ahmadi \cite{alalyan2020robust} proposed a URL-only CNN trained on over two million URLs.

Despite high accuracy, content-based techniques exhibit significant practical liabilities: (i) they require live connections to hostile servers, risking malware execution; (ii) cloaking mechanisms serve benign decoys to crawlers; and (iii) fetching and rendering pages introduces heavy network and computational latency, preventing inline inspection.

\subsection{URL and Network-Centric Machine Learning}
To achieve high-throughput triage without page rendering, researchers turned toward feature engineering on URLs and network metadata. Sahoo et al. \cite{sahoo2017malicious} and Rao and Pais \cite{rao2019detection} established taxonomies of URL lexical properties. Sahingoz et al. \cite{sahingoz2019machine} evaluated tree classifiers and linear SVMs on lexical representations. Babagoli et al. \cite{babagoli2019heuristic} proposed heuristic non-linear regression strategies for phishing detection, while Chiew et al. \cite{chiew2019new} analyzed hybrid ensemble feature selection. Vrban{\v{c}}i{\v{c}} et al. \cite{vrbancic2020datasets} standardized a benchmark containing 111 structured features across 58,645 URLs, capturing URL, domain, directory, file, query, and resolver metadata without page rendering.

\subsection{Explainable Artificial Intelligence (XAI) in Cybersecurity}
Recent works have explored post-hoc interpretability in phishing detection. Al-Fayoumi et al. \cite{alfayoumi2024xaiphd} proposed \textit{XAI-PhD}, leveraging TreeSHAP with ensemble models to interpret URL classification decisions and prune features to reduce prediction overhead. Uddin et al. \cite{uddin2025explainable} combined gradient-boosting models with SHAP on a public Kaggle dataset. Calzarossa et al. \cite{calzarossa2024explainable} proposed an explainable machine learning procedure to identify critical phishing indicators, while Kehkashan et al. \cite{kehkashan2025explainable} evaluated random forest with SHAP on a Kaggle dataset of 11,000 URLs.

While these studies established the utility of post-hoc explanations, several methodological distinctions and gaps remain: (i) prior works in this domain, including XAI-PhD \cite{alfayoumi2024xaiphd}, do not explicitly specify whether global SHAP attributions and subsequent feature pruning are restricted strictly to an independent training partition or evaluated globally across the dataset prior to final validation, a practice that risks subtle selection leakage if test instances inform feature pruning; (ii) attribution ranking stability across random seeds and subsamples is often unquantified; (iii) existing feature selection methods treat all features uniformly, whereas in practice, synchronous external network lookups (DNS/WHOIS) incur tens to hundreds of milliseconds of latency compared to microsecond in-memory lexical parsing; and (iv) computational latency reporting frequently omits repeated measurement intervals or single-query benchmarking. This paper addresses these challenges by enforcing a strict train-only TreeSHAP protocol on the standardized 58,645 URL corpus \cite{vrbancic2020datasets}, verifying attribution stability across seeds, and evaluating an uncertainty-gated cascade that explicitly decouples in-memory lexical parsing from network resolver queries.

